\section{Related Work}\label{sec:related}
A post-quantum anonymous credential can already be built and proved; what no prior work asks is whether that proof can be generated on the holder's own consumer machine. On a service operator's server, a changed predicate is absorbed by rebuilding and re-auditing the static verification circuit, so flexibility buys little there. Only once the holder proves on hardware they own does the cost of every such change start to matter. Each regenerated circuit artefact must then be redistributed to and re-trusted by every holder's device, with no operator on hand, as there would be behind a service boundary, to re-audit and swap it. This is the case where a zkVM's ability to edit the program, rather than rebuild the circuit, matters, and it is why a zkVM is needed here. Table~\ref{tab:positioning} places some prior work against that baseline: every construction shares the post-quantum, privacy-preserving goal (two rows reach credential-adjacent objects rather than a full anonymous credential, marked $\circ$), so the table maps only the axes that separate them.

The \emph{flexibility} column is where this thesis stands alone, and the last column names what it currently stops short of: a single proof that is at once post-quantum-sound and zero-knowledge. The sections below take the attempts in turn, asking at each step why the attempt stops short of the object measured here: constructing a credential directly, proving a signature in a static circuit, running one inside a zkVM, accelerating that execution with a precompile, then what security a zkVM receipt carries, and finally the gap this thesis fills.

\begin{table}[H]
\centering\small \setlength{\tabcolsep}{5pt} \renewcommand{\arraystretch}{1.2}
\begin{tabular}{@{}p{2.5cm}>{\centering\arraybackslash}p{1.9cm}p{3cm}>{\centering\arraybackslash}p{1.5cm}|>{\centering\arraybackslash}p{1.9cm}@{}}
\hline
\textbf{Work} & \textbf{Anonymous credential?} & \textbf{Substrate} & \textbf{Flexibility} & \textbf{PQ zero-knowledge} \\ \hline
Jeudy et al.~\cite{jeudy2023lattice} & $\bullet$ & --- & Low & $\bullet$ \\
BDEC~\cite{10.1007/978-981-96-0957-4_3} & $\bullet$ & Aurora & Low & $\bullet$ \\
Policharla et al.~\cite{policharla2023pqprivacypass} & $\circ$ & STARK & Low & $\bullet$ \\
CAPSS~\cite{feneuil2025capss} & $\circ$ & SmallWood & Low & $\bullet$ \\ \hline
\textbf{This thesis} & $\bullet$ & \textbf{zkVM + precompile} & \textbf{High} & $\oslash$ \\ \hline
\end{tabular}
\caption{Prior work mapped against the axes that separate it from this thesis; every row is a post-quantum construction, so that shared baseline is stated here. Symbols: $\bullet$ = provided; $\circ$ = a credential-adjacent object rather than a full anonymous credential; $\oslash$ = not achieved (Section~\ref{sec:security}); --- = not applicable. \label{tab:positioning}}
\end{table}

\subsection{Constructing an Anonymous Credential Directly}
The oldest approach builds an anonymous credential as a dedicated protocol. Chaum introduced the idea~\cite{chaum1985}, Camenisch and Lysyanskaya formalised it~\cite{camenisch2004signatures}, and the line runs through BBS, Pointcheval--Sanders, and Coconut~\cite{bbs04,pointchevalsanders2016,sonnino2019coconut} into the W3C Verifiable Credentials data model~\cite{w3cvc}; Rosenberg et al.~\cite{rosenberg2023zkcreds} reach the same goal by proving statements over existing identity infrastructure. All of these rest on classical assumptions, RSA, pairing, or discrete logarithm, and so do not survive a quantum adversary, which is why we treat them as deployment baselines rather than post-quantum baselines. The post-quantum case has been pursued largely with lattices: Jeudy et al.\ at Crypto~2023~\cite{jeudy2023lattice} give an efficient scheme on SIS and LWE, refined by later constructions~\cite{BootleLNS23,ArgoGJLRS24}. These deliver the credential, but at large proof and key sizes, and each is a bespoke protocol rather than an existing scheme's verification run as a program.

\subsection{Proving Signature Verification in a Static Circuit}
A second approach reaches a post-quantum credential generically, by proving in zero knowledge that one knows a valid signature, and compiles that verification into a fixed arithmetic circuit. The standardised post-quantum signatures ML-DSA, SLH-DSA, and Falcon~\cite{fips204,fips205,fips206} verify cheaply in the clear but embed into such a circuit only at large constraint counts, so this line instead uses signatures written for a low constraint count: the MPC-in-the-head family (Picnic, Banquet, Limbo~\cite{chase2017picnic,baum2021banquet,desaintguilhem2021limbo}, after Katz--Kolesnikov--Wang~\cite{katz2018kkw}) and the Legendre/power-residue family begun by LegRoast~\cite{beullens2020legroast}. Loquat~\cite{cryptoeprint:2024/868} bases unforgeability on the Legendre PRF and is engineered to fit a small rank-1 constraint system (R1CS); PLUM~\cite{10.1007/978-981-95-2961-2_6}, the scheme this thesis deploys, replaces that PRF with the $t$-th power-residue PRF and FRI with STIR~\cite{cryptoeprint:2024/390} (two protocols for the same task, testing that a committed function is close to a low-degree polynomial; Section~\ref{sec:prelim}), reporting at 128-bit security signatures roughly $1.5\times$ smaller and $1.28\times$ fewer constraints, with verification estimated up to $4\times$ faster, and the family has since been extended to ring signatures~\cite{cryptoeprint:2025/1841}.

Built on such a signature, several works reach a credential or credential-adjacent application. BDEC (Li, Zhang et al.)~\cite{10.1007/978-981-96-0957-4_3} is the closest prior work. It instantiates a post-quantum anonymous credential with Loquat and the transparent argument Aurora~\cite{cryptoeprint:2018/828}, proving unforgeability, anonymity, unlinkability, conditional linkability, and revocability, improving on the lattice baseline. But BDEC's guarantees hold for a static circuit, so any change to the signature or the disclosed predicate recompiles and re-audits it, and its anonymity requires the argument to be zero-knowledge, which the consumer-hardware zkVM receipt this thesis produces does not yet satisfy (Section~\ref{sec:security}). CAPSS (Feneuil and Rivain)~\cite{feneuil2025capss} pushes the constraint count $5$--$8\times$ below Loquat. It reports measured figures on SmallWood, its authors' own hash-based proving system, for an application built on the signature rather than a full credential scheme. So measured proving of an algebraic-hash signature (one built from a hash defined by algebraic operations, in the vocabulary of Section~\ref{sec:prelim}) predates this thesis. But it uses a different signature family from the PLUM scheme this thesis deploys, optimised for static-circuit constraint count. Policharla et al.~\cite{policharla2023pqprivacypass} take the same static route to a Privacy Pass, credential-adjacent rather than a full anonymous credential, proving a \emph{modified} Dilithium inside a hand-rolled STARK built on the Winterfell prover library. Across this line the verification is proved as a native circuit and recompiled on every change (preprocessing variants such as Fractal~\cite{cryptoeprint:2019/1076} and Marlin~\cite{chiesa2020marlin} only trade that recompile for a large per-relation setup).

\subsection{Running a Signature Inside a zkVM}
A third approach runs a post-quantum signature as an ordinary program inside a zkVM. s2morrow~\cite{s2morrow2025} verifies Falcon-512 and SPHINCS+ in Cairo with no custom precompile, HAPPIER~\cite{saygan2026happier} verifies XMSS through a built-in SHA-256 precompile, and an ML-DSA implementation~\cite{kota2025dilithiumzk} runs on SP1~\cite{succinct2024sp1} with its number-theoretic transform (NTT) implemented as guest code. These efforts run on Cairo, RISC~Zero~\cite{bruestle2023risczero}, and SP1. The broader zkVM landscape also includes Jolt~\cite{cryptoeprint:2023/1217} and the custom-chip-native OpenVM~\cite{openvm2025}, while compiler front-ends such as zkLLVM~\cite{Kaskov2024zkLLVM}, Circom~\cite{10002421}, and Noir~\cite{noir_lang} lock the predicate at compile time and so behave like the static-circuit regime above. Yet each demonstration verifies a \emph{standardised or hash-based} signature, authors no custom precompile for a large-prime-field algebraic hash, and is a bare signature rather than an anonymous credential.

\subsection{Accelerating a zkVM with a Precompile}
Where a primitive dominates a zkVM's cost, prior work replaces its slow emulated version with a custom accelerator (a precompile) that the program calls directly: SP1 ships SHA-256 and ECDSA precompiles~\cite{succinct2024sp1}, and a community RSA precompile is reported to cut verification by roughly two orders of magnitude.\footnote{Reported in vendor and community engineering write-ups without a peer-reviewed measurement; we cite it only as a qualitative indication of the pattern.} Such a precompile joins the trusted computing base and must be audited: Arguzz~\cite{hochrainer2025arguzz} found eleven soundness or completeness bugs across three of the six zkVMs it tested, which is why each precompile of this thesis ships a written soundness argument. Two 2025 lines automate the accelerate-or-not decision, powdr's autoprecompiles~\cite{powdr2025autoprecompiles} and the vApps mispricing analysis~\cite{vapps2025}; both rank or reprice from profiling data gathered \emph{after} the accelerator exists, whereas the criterion of Section~\ref{sec:theoretical} decides from the specification, before any build, whether a precompile over the signature's own large field removes the field-mismatch tax at all. No prior precompile targets a $199$-bit-prime algebraic hash, and none carries a from-specification pay criterion.

\subsection{The Security a zkVM Receipt Carries}\label{ssec:zkvm-proof-security}
A separate line asks not what a zkVM proves but what security its receipt actually carries. Hab\"ock and Kindi~\cite{cryptoeprint:2024/1037} show that the extra step these proof systems need to keep the witness hidden is easy to get wrong and a recurring source of implementation error, and RISC~Zero's own advisory~\cite{risc0zkadvisory} records a period in which receipts advertised as zero-knowledge leaked witness information; production toolchains therefore restore zero-knowledge with an outer pairing-based wrap in the Groth16 family~\cite{groth2016,succinct2025prooftypes}, whose security does not survive a quantum adversary, and only emerging hash-based zero-knowledge layers~\cite{cryptoeprint:2026/683} would remove that pairing. Each of these facts is established on its own; to our knowledge no prior work states their conjunction as an obstruction to a \emph{post-quantum} application, nor analyses which security properties of a credential survive execution on a zkVM, the question Section~\ref{sec:security} takes up. At the information-theoretic limit, everlasting privacy against a later unbounded adversary was introduced by Moran and Naor~\cite{morannaor2006} and shown unattainable from standard setup assumptions by M\"uller-Quade and Unruh~\cite{muellerquade2010}, with recent work bringing it to credential-adjacent tokens~\cite{cryptoeprint:2025/1030}; the everlasting-anonymity non-establishability result of Section~\ref{sec:security} is the analogue of that boundary for hash-committed credential showings, which no prior work states.

\subsection{The Gap and the Unreported Dimension}\label{ssec:gap}
The prior work has thus reached every piece of the problem and none of the whole: credentials built directly but never run in a zkVM; signature verification proved efficiently but in a recompiled static circuit; signatures run in a zkVM but standardised, without a custom precompile, and not as credentials; precompiles built but never for a large-field algebraic hash; and receipt security studied but never as a credential-anonymity obstruction. No prior work runs an algebraic-hash post-quantum anonymous credential inside a zkVM to learn whether it runs at all, at what proving cost, and with which security properties intact. A further dimension goes unreported throughout: the cost each change to the verification predicate imposes, which crypto-agility work~\cite{ott2019cryptoagility} names as a design concern but no evaluation measures, and which the recent zkVM systematisation~\cite{yang2026sokzkvm} omits alongside its classical-workload benchmarks. This thesis fills the gap and reports that dimension as \emph{update churn}: it develops a cost criterion for when a precompile pays (Section~\ref{sec:theoretical}), analyses which security properties survive the move (Section~\ref{sec:security}), and measures the deployability and the update churn of BDEC-with-PLUM on consumer hardware (Section~\ref{sec:evaluation}).
